\section{Introduction}
\label{Introduction}
%The introduction comprises mainly two main parts: (1) The motivation answers the question ”Why is the problem a problem?”. I.e. it outlines why the content of the paper is relevant, both from an academic point of view but also from a market point of view. Normally this requires quotes and references from studies written by well-know researchers, associations, organization and politics since in most cases the author does not have the standing to make such claims. (2) The second part defines the research questions or the research hypotheses.

The traditional design of new materials and their production are based on exploratory research and heuristic methodology \cite{morgan2020opportunities,butler2018machine,MolMatSurv}. Newer approaches use computer-aided modeling of materials properties and integrate them into the materials design process. This enables a more efficient screening process for possible theoretical solutions \cite{ramprasad2017machine, sousa2021material}. But materials properties are still described spatially using physical models. These models are always based on an idealized theoretical morphology. 

This is where state-of-the-art laboratory-based computed tomography (CT) can create CTs of the real morphology. It is an accurate and fast way to digitize materials and assemblies morphologies over statistically relevant volumes and with appropriate resolution \cite{tom3, tom2, tom1}. Once a digital image of the morphology is available, superimposed numerical solutions of the partial differential equations can approximate the effective materials properties \cite{num2, num1}. However, numerical solving is associated with an enormous computational effort. In addition, uncertainties in the theoretical description with physical models, for example due to the system complexity or the non-ideal behavior of some components, remain undescribed.

Data-driven methods close the gap between empirical observation and physical modelling. Neural networks have become the dominant approach for data-based modeling in many complex areas \cite{YOLO,UNET, AlphaGo}. Their success is largely due to their ability as universal function approximators and the existence of highly efficient optimization algorithms \cite{Optimizers}. However, in many neural network architectures, the purely data-based approach limits their ability to extrapolate to data that is very different from the training data \cite{IPNN}. Nevertheless, neural networks are successfully used to create models for material properties \cite{MLMF2,MLMF1, MatMLSurv,MolMatSurv}.

Unfortunately, the black-box nature of neural networks means that the learned function is usually uninterpretable by a human. This makes it difficult to use the trained model to gain deeper insights into the physical principles that lead to the observed properties of the materials, or even to assess the reliability of the model prediction. As a response to the aforementioned shortcomings of standard neural networks, several alternative approaches have been developed that improve these models by learning physical concepts of varying degrees of interpretability. \cite{PIM, GAN, Graph}

One promising method in this domain is representation learning (RepL). The aim of this field is to create models which are capable of mapping observations of a system into representations that are useful for making downstream predictions for the system \cite{RepL}. These representations usually offer an improved degree of interpretability and are often explicitly designed for that purpose \cite{kingma:2014a, BVAE, Eureqa, AIFeynman}. A number of research hypotheses (RH) arise from this context. The proposed solution takes these into account and enables a faster and more efficient material design process.

\textit{RH 1:} The solution is based on a formalization of materials designs that capture the relevant materials properties and serve as input to a surrogate model. Such a model must be developed based on requirements and uncertainties from manufacturing processes.

\textit{RH 2:} The network typologies (Bayesian Generative Autoencoders) and the corresponding algorithms (active learning, RepL) must be developed for this domain. Neural networks for time series -- which would be suitable -- are an on-going research field in ML \cite{HEWAMALAGE2021388,8489399}. Here too, a training strategy must be developed.

\textit{RH 3:} To extract an explicit, formula-based material model from the neural network, several research directions exist: From symbol-regression based-approaches \cite{doi:10.1073/pnas.1517384113}, via questions-based algorithms \cite{PhysRevLett.124.010508} to approaches which define priors on latent variables \cite{kingma:2014a}. Here a suitable algorithm and neural network type must be identified and verified.

\textit{RH 4:} Materials morphologies need to be encoded into a set of feature vectors that are automatically extracted from CT data. Some of these features are straight-forward (e.g., density), while others will require novel image-processing strategies. In related applications, convolutional neural networks have proven to be particularly suitable for this.

\textit{RH 5:} CT data can be used to quantify the effects of manufacturing processes. This knowledge enables high-resolution conditional generative neural networks (generative adversarial networks (GANs) \cite{goodfellow2014generative} or denoising diffusion probabilistic models (DDPM) \cite{ho2020denoising}) to synthesize CTs based on material and manufacturing parameters. For this purpose, a suitable experimental design must be created. In addition, the algorithms must consider highly modal and diverse conditions.
